\subsection{从固定阶失效到 Sudakov 重求和}
\label{sec:qcd-sudakov-thrust-applications}

软、共线极限已说明，降低喷注分辨阈值会产生大对数。仅仅指出固定阶展开失效还不能给出可用预测；下一步必须把各阶重复出现的辐射贡献重新组织。一个适合完整演算的观测量是推力（thrust）：它测量强子末态的动量沿某条最优轴集中到什么程度。

对质心系中所有测得的末态三动量 $\mathbf p_i$，定义
\begin{equation}
 T_{\rm thr}=\max_{|\mathbf n|=1}
 \frac{\sum_i|\mathbf p_i\cdot\mathbf n|}{\sum_i|\mathbf p_i|},
 \qquad \tau_T=1-T_{\rm thr}.
 \label{eq:qcd-thrust-definition}
\end{equation}
最大化得到的 $\mathbf n$ 称为推力轴。无质量末态有 $1/2\leq T_{\rm thr}\leq1$；两个背对背的窄喷注对应 $\tau_T\to0$。增加一个无限软粒子不改变该比值，共线分裂也不改变投影总和，因此推力是红外--共线安全的。有限质量强子的实验定义仍使用实际三动量，不能在精密比较中不加说明地把能量与动量模长互换。

累积分布的精确定义为
\begin{equation}
 \Sigma_T(\tau)=\frac{1}{\sigma_{\rm had}}
 \sum_X\int\dd\sigma_X\,
 \Theta[\tau-\tau_T(X)].
 \label{eq:qcd-thrust-cumulant}
\end{equation}
$\sigma_{\rm had}$ 是同一初态下的包容强子截面，$\dd\sigma_X$ 已含振幅平方、通量与相空间；$\Theta$ 是阶跃函数。这个定义对多粒子末态直接测量 $\tau_T(X)$，尚未假设只辐射一个胶子。精确分布非负、随 $\tau$ 单调增加，并在覆盖全部物理末态时趋于一。

在 $\tau\ll1$ 的两喷注区域，领先幂动力学分别含硬能标 $E_{\rm cm}$、喷注能标 $E_{\rm cm}\sqrt\tau$ 和软能标 $E_{\rm cm}\tau$。在这些尺度都处于微扰区时，可以用硬函数、喷注函数和软函数作系统因子化；它们的重整化群演化组织完整的对数级数\cite{BecherSchwartz2008Thrust}。下面先保留固定 $\alpha_s$ 的领先双对数，逐项算出这一结构最直接的辐射解释。该精度保留 $\alpha_s^n\ln^{2n}\tau$，不包含完整的运行耦合领先对数或次领先对数。

对于相对某条硬夸克方向的小角软胶子，定义能量分数 $z_g=2E_g/E_{\rm cm}$、角变量 $u_\theta=\theta^2/4$。单次发射对推力亏损的领先贡献为 $v=z_gu_\theta$，单条夸克腿的实辐射测度为
\begin{equation}
 \dd\mathcal P_q=\frac{\alpha_sC_F}{\pi}
 \frac{\dd z_g}{z_g}\frac{\dd u_\theta}{u_\theta}.
 \label{eq:qcd-thrust-emission-measure}
\end{equation}
有限角度与非软分裂项会产生较低次对数。只求双对数时，可以把两个变量的上限设为一；改变这些数量级为一的上限会改变单对数与常数项，不能用于确定更高精度的系数。

两条背对背的夸克腿都能辐射。若要求没有单次发射使 $v>\tau$，排除区域中的总辐射权重为
\begin{equation}
 R_T(\tau)=\frac{2\alpha_sC_F}{\pi}
 \int_\tau^1\frac{\dd z_g}{z_g}
 \int_{\tau/z_g}^1\frac{\dd u_\theta}{u_\theta}
 =\frac{\alpha_sC_F}{\pi}\mathcal L_T^2,
 \qquad \mathcal L_T=\ln\frac1\tau.
 \label{eq:qcd-thrust-radiator}
\end{equation}
虚修正和任意多个允许的未分辨实发射共同产生 Sudakov 因子，固定耦合双对数结果为
\begin{equation}
 \Sigma_T^{\rm DL}(\tau)=\ee^{-R_T(\tau)},\qquad
 \frac{\dd\Sigma_T^{\rm DL}}{\dd\tau}
 =\frac{2\alpha_sC_F}{\pi}\frac{\mathcal L_T}{\tau}
 \ee^{-(\alpha_sC_F/\pi)\mathcal L_T^2}.
 \label{eq:qcd-thrust-sudakov}
\end{equation}
在此精度内可以将包容归一化替换成 Born 归一化，因为两者差异从不含双对数的相对 $\order(\alpha_s)$ 开始。重求和并没有消除胶子发射，而是计算观测窗口中始终没有出现可分辨发射的概率。

固定阶展开只保留 $\Sigma_T=1-R_T+\cdots$；当 $R_T>1$ 时会给负的累计率。指数把整组 $R_T^n/n!$ 保留下来，使固定耦合近似的累计率保持在零与一之间。然而概率非负只是一项必要检查，不能据此认定该近似已达到实验精度。在 $\tau\to0$ 时，真实运行耦合会取样低能区域；当 $E_{\rm cm}\tau$ 接近强子化能标，微扰 Sudakov 公式也需要非微扰输入。

\begin{exbox}[大对数怎样使一阶近似失效]
取固定教学耦合 $\alpha_s=0.120$、$C_F=4/3$。在 $\tau=0.10$ 时，$R_T\simeq0.270$，一阶累计率为 $0.730$，双对数指数结果为 $0.763$；在 $\tau=0.01$ 时，$R_T\simeq1.080$，一阶结果变为 $-0.080$，指数结果仍为 $0.340$。

固定阶失效的边界由 $R_T\sim1$ 给出，
\begin{equation*}
 \tau\sim\exp\!\left[-\sqrt{\frac{\pi}{\alpha_sC_F}}\right]\simeq0.0119.
\end{equation*}
这解释了小耦合与固定阶可靠性之间为何还隔着一个分辨率条件。若 $E_{\rm cm}=\SI{100}{\giga\electronvolt}$，$\tau=0.01$ 对应的软能标仅为 $\SI{1}{\giga\electronvolt}$；即使硬尺度上的 $\alpha_s$ 很小，也不能据此忽略软区运行与强子化。
\end{exbox}

\begin{figure}[htbp]
 \centering
 \includegraphics[width=0.73\textwidth]{qcd_thrust_sudakov.pdf}
 \caption{固定耦合双对数近似下的推力累计率。虚线是一阶展开，实线是同一双对数级数的指数求和；负的一阶结果标志其失效。横轴只展示两喷注区域，曲线未包含运行耦合、次领先对数与强子化。}
 \label{fig:qcd-thrust-sudakov}
\end{figure}
\figref{fig:qcd-thrust-sudakov} 比较的是同一对数精度的截断与求和。它适合判断截断何时失去层级，定量拟合推力分布还需要更高阶因子化与匹配。

\begin{derivation}[双对数的三角形积分区域]
对于无质量末态，$\sum_i|\mathbf p_i|=E_{\rm cm}/c$。软胶子相对最近推力轴方向的贡献为
\begin{equation*}
 v=\frac{E_g}{E_{\rm cm}}(1-|\cos\theta|)
 \simeq\frac{E_g\theta^2}{2E_{\rm cm}}=z_gu_\theta.
\end{equation*}
从夸克分裂核的软端点 $P_{qq}\simeq2C_F/(1-z)$ 出发，使用胶子分数 $z_g=1-z$，得到\eqnrefs{eq:qcd-thrust-emission-measure}。对两条腿相加给因子二；宽角软干涉与推力轴反冲对当前双对数的系数没有额外贡献。

令 $r_z=\ln(1/z_g)$、$r_\theta=\ln(1/u_\theta)$。可分辨发射条件 $z_gu_\theta>\tau$ 变成
\begin{equation*}
 r_z\geq0,\qquad r_\theta\geq0,\qquad r_z+r_\theta<\mathcal L_T.
\end{equation*}
该区域面积为 $\mathcal L_T^2/2$。因此
\begin{equation*}
 R_T=\frac{2\alpha_sC_F}{\pi}\frac{\mathcal L_T^2}{2}
 =\frac{\alpha_sC_F}{\pi}\mathcal L_T^2.
\end{equation*}
面积中的 $1/2$ 与两条辐射腿的因子二恰好抵消；二者来源独立，遗漏任意一个都会使指数相差一倍。
\end{derivation}

\begin{derivation}[实虚抵消、指数化与多发射约束]
先引入有限的软、共线分辨下限，使全发射区域 $\Omega$ 的积分 $I_\Omega=\int_\Omega\dd\mathcal P$ 有限。在领先独立发射近似中，虚修正提供 $\ee^{-I_\Omega}$，$n$ 个相同胶子的实发射测度带 $1/n!$。若仅施加逐个发射的 $v_i<\tau$ 约束，累加任意发射数得到
\begin{align*}
 \Sigma_{\max}(\tau)
 &=\ee^{-I_\Omega}\sum_{n=0}^\infty\frac1{n!}
 \left[\int_\Omega\dd\mathcal P\,\Theta(\tau-v)\right]^n\\
 &=\exp\!\left[-\int_\Omega\dd\mathcal P\,\Theta(v-\tau)\right].
\end{align*}
移除辅助红外下限以后，只留下有限的排除区积分 $R_T$。这个实虚相消使结果不依赖“总共发射了多少个无限软胶子”。

实际推力的领先软贡献是 $\tau_T=\sum_i v_i$，不是 $\max_i v_i$。对于强有序的双对数主导区域，最大项决定总和，所以 $\Sigma_{\max}$ 给出\eqnrefs{eq:qcd-thrust-sudakov} 的精度。若要保留多个同量级发射对总和的影响，应使用拉普拉斯变量 $\nu_L$：
\begin{align*}
 \mathcal Z_T(\nu_L)
 &=\ee^{-I_\Omega}\sum_{n=0}^\infty\frac1{n!}
 \prod_{i=1}^n\int_\Omega\dd\mathcal P_i\,
 \ee^{-\nu_L\sum_i v_i}\\
 &=\exp\!\left[\int_\Omega\dd\mathcal P\,(\ee^{-\nu_Lv}-1)\right],\\
 \Sigma_T(\tau)
 &=\int_{a_{\rm inv}-\ii\infty}^{a_{\rm inv}+\ii\infty}
 \frac{\dd\nu_L}{2\pi\ii\nu_L}\,
 \ee^{\nu_L\tau}\mathcal Z_T(\nu_L),\qquad a_{\rm inv}>0.
\end{align*}
这里 $a_{\rm inv}$ 为无量纲拉普拉斯逆变换路径的正实部。$\mathcal Z_T$ 在独立软发射模型内保留了加和约束，但尚不包含 QCD 的全部相关多胶子发射。系统提高对数精度需要同时处理运行耦合、非软分裂、软相关及相应硬、喷注、软函数。
\end{derivation}

\begin{derivation}[固定阶匹配与强子化误差的尺度]
重求和适用于小 $\tau$，完整固定阶结果则保留远离两喷注极限的非奇异项。若 $\Sigma_T^{\rm FO,(n)}$ 已知到 $\alpha_s^n$，一个加性匹配定义是
\begin{equation}
 \Sigma_T^{\rm match}
 =\Sigma_T^{\rm resum}+\Sigma_T^{\rm FO,(n)}
 -\left.\Sigma_T^{\rm resum}\right|_{\alpha_s^0,\ldots,\alpha_s^n}.
 \label{eq:qcd-thrust-matching}
\end{equation}
把重求和结果展开为 $\sum_{k\geq0}\alpha_s^kr_k$，最后一项逐项去掉 $k\leq n$ 的重叠贡献，故
\begin{equation*}
 \left.\Sigma_T^{\rm match}\right|_{k\leq n}
 =\Sigma_T^{\rm FO,(n)},\qquad
 \left.\Sigma_T^{\rm match}\right|_{k>n}
 =\sum_{k>n}\alpha_s^kr_k.
\end{equation*}
这同时说明不能直接把固定阶曲线与重求和曲线相加：共有的软、共线项会被重复计数。匹配时须使用相同的截面归一化、耦合方案和尺度定义。

若非微扰软动量引起的推力位移为 $\delta\tau\sim\Lambda_E/E_{\rm cm}$，其中 $\Lambda_E$ 是强子化能标，则在宽度为 $\tau$ 的两喷注尾区，其相对影响由
\begin{equation}
 \epsilon_{\rm had}\sim\frac{\Lambda_E}{E_{\rm cm}\tau}
 \label{eq:qcd-thrust-hadronization-control}
\end{equation}
控制。即使总截面的幂修正很小，小 $\tau$ 微分分布也会放大低能效应。$\epsilon_{\rm had}\ll1$ 时可以组织矩展开；接近一时必须保留非微扰软函数的形状。这是“硬过程可微扰”与“所选观测量可微扰”之间的具体尺度区别。
\end{derivation}
